\section{The algorithm and its guarantee}
\label{d:sec-algorithm}

We now assemble the preceding constructions into a deterministic algorithm.
The input consists of a vertex set \(V=\{1,\ldots,n\}\) and nonnegative
integer multiplicities \(k_{\{i,j\}}\) for distinct vertex pairs, written in
binary; an omitted pair has multiplicity zero.  Listing parallel edges
individually is also allowed, since their multiplicities can first be
counted.  Let \(N\) be the number of perfect matchings, with parallel
edge choices distinguished.  Fix the vertex and pair orderings, and use
the deterministic cut procedure of Section~\ref{d:sec-odd-cuts}.

\begin{enumerate}[label=\textbf{\arabic*.},leftmargin=*]
\item \textbf{Handle the elementary cases.}
If \(n=0\), return \(1\).  If \(n\) is odd, return \(0\).
If \(n=2\), return \(k_{\{1,2\}}\).
Sort the endpoint labels occurring in positive-multiplicity pair records
and count the distinct labels. If fewer than \(n\) occur, an isolated
vertex forces the count to be zero; return \(0\).
This test uses the input records even when \(n\) is supplied in binary,
before allocating any array indexed by all vertices or pairs.

\item \textbf{Set up the complete-pair objective.}
Now \(n\ge4\) is even.  Use all \(p=\binom n2\) pair coordinates.
Compute \(b_{\max},L,\beta\) from Equation~\eqref{d:eq-completion-weights},
and let \(K\) be the least nonnegative integer with \(2^K\ge p\).
Define \(\psi\) and \(f\) by Equations~\eqref{d:eq-dyadic-chords}
and~\eqref{d:eq-surrogate-objective}.
Set \(A_0=10n^2\), define \(w\) by Equation~\eqref{d:eq-max-violation},
and compute \(M_0,\Phi,D_0\) from Equation~\eqref{d:eq-penalty-parameters}.
Set
\[
                 N_{\mathrm{it}}=pD_0,\qquad a^{(0)}=0\in\mathbb Z^p.
\]
The integer vector \(a^{(j)}\) represents \(x^{(j)}=a^{(j)}/D_0\).

\item \textbf{Evaluate and update \(N_{\mathrm{it}}\) trial points.}
For \(j=0,\ldots,N_{\mathrm{it}}-1\), do the following.
\begin{enumerate}[label=(\alph*),leftmargin=*]
\item Apply Lemma~\ref{d:lem-odd-cut-oracle} to the integral capacities
\(a^{(j)}_e\), with every vertex terminal, obtaining a minimum-capacity
odd cut \(S_j\).  At \(x=x^{(j)}\), compare
\[
 0,\quad x(\delta(1))-1,\quad 1-x(\delta(1)),\quad\ldots,\quad
 x(\delta(n))-1,\quad 1-x(\delta(n)),\quad 1-x(\delta(S_j)).
\]
Their maximum is \(w(x)\).  Select the first maximizing expression in
this list and let \(r^{(j)}\in\{-1,0,1\}^p\) be its linear coefficient
vector.  Evaluate \(\Phi(x)\) exactly and retain the largest trial value
seen so far.
\item In each coordinate choose the right-hand slope \(\sigma_e^{(j)}\)
of \(\psi\) at \(x_e^{(j)}\), using the left-hand slope at \(1\).
Set
\[
             g_e^{(j)}
                =\sigma_e^{(j)}+\beta_e-M_0r_e^{(j)}.
\]
These are the integral supergradients from
Lemma~\ref{d:lem-exact-penalty}.
\item Update the integer numerators by
\[
 a_e^{(j+1)}
       =\min\!\left\{D_0,\,
             \max\{0,a_e^{(j)}+g_e^{(j)}\}\right\}
       \qquad(e\in E(K_n)).
\]
\end{enumerate}

\item \textbf{Return the estimate.}
Let
\[
           R=\left\lfloor
                  \max_{0\le j<N_{\mathrm{it}}}\Phi(x^{(j)})
              \right\rfloor.
\]
If \(R<-10n\), return \(0\).  Otherwise return
\[
                      A=2^{\rho},
             \qquad \rho=\max\{0,R-5n\}.
\]
\end{enumerate}

The slope convention and the ordering in Step~3 settle all remaining
ties.  The numerator update is exactly the clipped step
\(x^{(j)}+g^{(j)}/D_0\) of Lemma~\ref{d:lem-projected-ascent}.

\begin{proof}[Approximation guarantee in Theorem~\ref{thm:algorithm-guarantee}]
The elementary cases return the exact count, so suppose \(n\ge4\) is
even and the algorithm reaches the optimization.
Lemma~\ref{d:lem-projected-ascent}, Lemma~\ref{d:lem-dyadic-entropy}, and
Corollary~\ref{d:coarse-weighted} give, respectively,
\[
 R\le f^*\le R+2,\qquad
 f^*\le q^*\le f^*+n,\qquad
 q^*-4n\le\log_2 Z_\beta\le q^*+n.
\]
Combining them yields
\begin{equation}
\label{d:eq-algorithm-partition-bracket}
       R-4n\le\log_2 Z_\beta
              \le R+2n+2\le R+3n.
\end{equation}
If \(N=0\), Lemma~\ref{d:lem-penalized-count} gives
\(\log_2 Z_\beta\le-20n\), hence \(R\le-16n<-10n\).
If \(N>0\), the same lemma gives \(Z_\beta\ge1\), hence
\(R\ge-3n>-10n\).  Thus the zero test is exact.

Now let \(N>0\).
The inequalities \(Z_\beta\le2N\) and
\(Z_\beta\ge2^{-n/2}N\), again from
Lemma~\ref{d:lem-penalized-count}, imply
\begin{equation}
\label{d:eq-algorithm-count-bracket}
 R-5n
   \le R-4n-1
   \le\log_2 N
   \le R+\frac{7n}{2}
   \le R+4n.
\end{equation}
Also \(\log_2N\ge0\), so
\(\rho=\max\{0,R-5n\}\le\log_2N\).
Since \(\rho\ge R-5n\), Equation~\eqref{d:eq-algorithm-count-bracket} gives
\(\log_2N-\rho\le9n\).
Exponentiating proves \(A\le N\le2^{9n}A=512^nA\).
\end{proof}

\begin{proposition}[Bit complexity]
\label{d:prop-bit-complexity}
The algorithm above is uniform and deterministic, and its running time
and output length are polynomial in the input bit length.
\end{proposition}

\begin{proof}
The elementary cases require only integer operations on the input.
In the remaining case put \(\tau=n+b_{\max}+1\), so \(\tau\ge5\).
The parameters have polynomial numerical magnitude, not just polynomial
bit length.  Indeed \(p\le\tau^2\), \(K\le n\le\tau\),
\(2^K<2p\), \(L=100\tau^3\), and \(A_0\le10\tau^2\).
Consequently
\begin{equation}
\label{d:eq-parameter-size-bounds}
\begin{aligned}
 M_0
   &=p(K+L)(2+A_0)
        \le1212\,\tau^7,\\
 D_0
   &=p(K+L+M_0)^2
        \le1313^2\,\tau^{16},\\
 N_{\mathrm{it}}
   &=pD_0
        \le1313^2\,\tau^{18}.
\end{aligned}
\end{equation}
Here \(K+L\le101\tau^3\) and \(2+A_0\le12\tau^2\) give the first
inequality, and \(K+L+M_0\le1313\tau^7\) gives the second.
The exponent of each positive multiplicity is obtained from its binary
length: \(\lfloor\log_2 k_e\rfloor=\operatorname{bitlen}(k_e)-1\).
No logarithm is evaluated numerically.

The update preserves
\(a_e^{(j)}\in\{0,\ldots,D_0\}\), because every supergradient coordinate is
integral.  Thus every iterate has denominator dividing \(D_0\).
On a noninitial dyadic interval the affine formula for \(\psi\) is
\((k-1)t+2^{-k}\), and on the interval next to zero it is \(Kt\).
The intercept denominators therefore divide \(2^K\).
Every degree or cut sum, and hence \(w(x^{(j)})\), has denominator
dividing \(D_0\).  A single common denominator for all trial objectives is
\begin{equation}
\label{d:eq-fixed-objective-denominator}
                             Q=D_0\,2^K.
\end{equation}
This denominator does not grow with the number of iterations or with
changes of dyadic piece.  The piece comparisons, objective comparisons,
and the final floor can all be made by exact integer arithmetic.

For an explicit bound on the operands, \(0\le w(x)\le n\) on the box,
and
\[
        |f(x)|\le p(1+L),\qquad
        |\Phi(x)|\le p(1+L)+nM_0,
        \qquad x\in[0,1]^p.
\]
Together with Equation~\eqref{d:eq-parameter-size-bounds} and
\(2^K<2p\), these bounds show that all objective numerators with
denominator \(Q\) have polynomial magnitude in \(\tau\).
The supergradients, scaled capacities, and counters do as well.
Their bit lengths are therefore \(O(\log\tau)\).

At any trial point, the cut oracle receives capacities \(a_e^{(j)}\),
not the original multiplicities.  Their total is at most \(pD_0\).
After any sequence of contractions, each surviving capacity is a sum
of original capacities between two disjoint vertex blocks; internal
edges are discarded.  Every original edge contributes to at most one
surviving pair, so the same total bound holds throughout the recursion.
By Lemma~\ref{d:lem-unit-mincut}, each terminal-pair computation uses at
most \(pD_0\) unit augmentations.  Lemma~\ref{d:lem-odd-cut-oracle} uses
\(O(n)\) recursive calls, each with \(O(n^2)\) terminal-pair
computations.  Dense residual path searches take \(O(n^2)\) elementary
operations.  The entire optimization thus uses at most
\[
            O\!\left(N_{\mathrm{it}}\,n^5(pD_0+1)\right)
                         =O(\tau^{41})
\]
elementary graph and arithmetic operations, on integers of
\(O(\log\tau)\) bits.  Standard integer arithmetic gives a polynomial
bit bound.

It remains to bound the length of the explicit integer output.
Lemma~\ref{d:lem-projected-ascent} gives \(R\le f^*\).
Since \(\psi(t)\le -t\log_2t\le1\) on \([0,1]\),
\(\beta_e\le b_{\max}\), and \(\sum_e x_e=n/2\) on \(P(K_n)\),
we have \(f^*\le p+b_{\max}n/2\).  This upper bound is nonnegative, so the
clamped exponent satisfies
\begin{equation}
\label{d:eq-clamped-output-size}
 \rho=\max\{0,R-5n\}
       \le\max\{0,R\}
       \le\max\{0,f^*\}
       \le p+\frac{b_{\max}n}{2}.
\end{equation}
Writing \(2^\rho\) as a binary integer therefore takes
\(O(n^2+nb_{\max}+1)\) bits and operations.  The zero and exact elementary
outputs also have polynomial length.

For an explicit vertex list, both \(n\) and \(b_{\max}\) are bounded by the
input bit length.  The same conclusion holds when \(n\) is given in
binary and only positive pair records are listed. The coverage test in
Step~1 takes polynomial time in the input bit length. If it does not
return zero, every vertex appears in a record, so \(n\) is at most twice
the number of positive pair records. Thus the complete-pair allocation
and the optimization are polynomial in this encoding as well.
\end{proof}

Theorem~\ref{thm:algorithm-guarantee} and
Proposition~\ref{d:prop-bit-complexity} give both the approximation
guarantee and its uniform polynomial bit implementation.
If an upper estimate is preferred, \(512^nA\) lies between \(N\)
and \(512^nN\).
